% !TeX program = lualatex
% !TeX root = predavanje04.tex
% ============================================================
% Zajednički preamble za sva predavanja RWA kolegija
% Uključiti s: % ============================================================
% Zajednički preamble za sva predavanja RWA kolegija
% Uključiti s: % ============================================================
% Zajednički preamble za sva predavanja RWA kolegija
% Uključiti s: \input{../common/preamble.tex}
% ============================================================

\documentclass[aspectratio=149,11pt]{beamer}
\usetheme{CambridgeUS}

\usepackage{fontspec}              % radi samo sa XeLaTeX/LuaLaTeX
\setsansfont{Latin Modern Sans}    % sans-serif font
\setmonofont{Latin Modern Mono}    % monospace font za kod
\usepackage[croatian]{babel}
\usepackage{csquotes}              % \enquote{} → hrvatski navodnici „..."
\usepackage{lmodern}
\usepackage{booktabs}
\usepackage{listings}
\usepackage{xcolor}
\usepackage{hyperref}

% ---------- Autor i logo ----------
\author{mag.ing. Josip Torić (SIT UNIZD)}
\titlegraphic{\includegraphics[height=2.5cm]{../common/SIT_logo_tekst1.pdf}}

% ---------- Listings ----------
\definecolor{codebg}{RGB}{248,248,248}
\definecolor{codeframeclr}{RGB}{220,220,220}
\lstdefinelanguage{bash}{
  morekeywords={sudo,apt,brew,pip,pipx,uv,winget,python,git,flask,source,deactivate,code,gunicorn,docker,compose,heroku,render,railway,node,npm,pnpm,yarn,uvicorn,alembic,pytest,ruff,black,precommit},
  sensitive=false,
  morecomment=[l]{\#}
}
\lstdefinelanguage{python}{
  morekeywords={from,import,def,class,return,if,elif,else,for,in,while,try,except,raise,with,as,not,and,or,is,lambda,yield,True,False,None,await,async},
  sensitive=true,
  morecomment=[l]{\#},
  morestring=[b]',
  morestring=[b]"
}
\lstdefinelanguage{sql}{
  morekeywords={SELECT,FROM,WHERE,INSERT,INTO,VALUES,UPDATE,SET,DELETE,CREATE,TABLE,PRIMARY,KEY,FOREIGN,REFERENCES,UNIQUE,NOT,NULL,INDEX,ON,ALTER,DROP,CONSTRAINT},
  sensitive=false,
  morecomment=[l]{--}
}
\lstdefinelanguage{json}{
  morecomment=[l]{//},
  morestring=[b]",
  showstringspaces=false
}
\lstdefinelanguage{typescript}{
  morekeywords={import,export,from,const,let,var,function,return,if,else,for,of,in,while,switch,case,break,async,await,try,catch,finally,throw,new,class,this,extends,implements,interface,type,enum,readonly,public,private,protected,static,get,set,true,false,null,undefined,void,string,number,boolean,any,unknown,never,typeof,keyof,as},
  sensitive=true,
  morecomment=[l]{//},
  morecomment=[s]{/*}{*/},
  morestring=[b]',
  morestring=[b]",
  morestring=[b]`
}
\lstalias{javascript}{typescript}
\lstset{
  backgroundcolor=\color{codebg},
  basicstyle=\ttfamily\footnotesize,
  frame=single,
  rulecolor=\color{codeframeclr},
  frameround=tttt,
  breaklines=true,
  tabsize=2,
  showstringspaces=false,
  upquote=true,
  columns=fullflexible
}

% ---------- Helper macros ----------
\newcommand{\codefile}[1]{\texttt{#1}}
\newcommand{\kbd}[1]{\texttt{#1}}

% ---------- Code frame environment ----------
% Automatski dodaje fragile i shrink=5, koristi se umjesto \begin{frame}
% Primjer: \begin{codeframe}{Naslov} ... \end{codeframe}
% Primjer s opcijama: \begin{codeframe}[shrink=15]{Naslov} ... \end{codeframe}
\newenvironment{codeframe}[2][]
  {\begin{frame}[fragile,shrink=5,#1]{#2}}
  {\end{frame}}

% ============================================================

\documentclass[aspectratio=149,11pt]{beamer}
\usetheme{CambridgeUS}

\usepackage{fontspec}              % radi samo sa XeLaTeX/LuaLaTeX
\setsansfont{Latin Modern Sans}    % sans-serif font
\setmonofont{Latin Modern Mono}    % monospace font za kod
\usepackage[croatian]{babel}
\usepackage{csquotes}              % \enquote{} → hrvatski navodnici „..."
\usepackage{lmodern}
\usepackage{booktabs}
\usepackage{listings}
\usepackage{xcolor}
\usepackage{hyperref}

% ---------- Autor i logo ----------
\author{mag.ing. Josip Torić (SIT UNIZD)}
\titlegraphic{\includegraphics[height=2.5cm]{../common/SIT_logo_tekst1.pdf}}

% ---------- Listings ----------
\definecolor{codebg}{RGB}{248,248,248}
\definecolor{codeframeclr}{RGB}{220,220,220}
\lstdefinelanguage{bash}{
  morekeywords={sudo,apt,brew,pip,pipx,uv,winget,python,git,flask,source,deactivate,code,gunicorn,docker,compose,heroku,render,railway,node,npm,pnpm,yarn,uvicorn,alembic,pytest,ruff,black,precommit},
  sensitive=false,
  morecomment=[l]{\#}
}
\lstdefinelanguage{python}{
  morekeywords={from,import,def,class,return,if,elif,else,for,in,while,try,except,raise,with,as,not,and,or,is,lambda,yield,True,False,None,await,async},
  sensitive=true,
  morecomment=[l]{\#},
  morestring=[b]',
  morestring=[b]"
}
\lstdefinelanguage{sql}{
  morekeywords={SELECT,FROM,WHERE,INSERT,INTO,VALUES,UPDATE,SET,DELETE,CREATE,TABLE,PRIMARY,KEY,FOREIGN,REFERENCES,UNIQUE,NOT,NULL,INDEX,ON,ALTER,DROP,CONSTRAINT},
  sensitive=false,
  morecomment=[l]{--}
}
\lstdefinelanguage{json}{
  morecomment=[l]{//},
  morestring=[b]",
  showstringspaces=false
}
\lstdefinelanguage{typescript}{
  morekeywords={import,export,from,const,let,var,function,return,if,else,for,of,in,while,switch,case,break,async,await,try,catch,finally,throw,new,class,this,extends,implements,interface,type,enum,readonly,public,private,protected,static,get,set,true,false,null,undefined,void,string,number,boolean,any,unknown,never,typeof,keyof,as},
  sensitive=true,
  morecomment=[l]{//},
  morecomment=[s]{/*}{*/},
  morestring=[b]',
  morestring=[b]",
  morestring=[b]`
}
\lstalias{javascript}{typescript}
\lstset{
  backgroundcolor=\color{codebg},
  basicstyle=\ttfamily\footnotesize,
  frame=single,
  rulecolor=\color{codeframeclr},
  frameround=tttt,
  breaklines=true,
  tabsize=2,
  showstringspaces=false,
  upquote=true,
  columns=fullflexible
}

% ---------- Helper macros ----------
\newcommand{\codefile}[1]{\texttt{#1}}
\newcommand{\kbd}[1]{\texttt{#1}}

% ---------- Code frame environment ----------
% Automatski dodaje fragile i shrink=5, koristi se umjesto \begin{frame}
% Primjer: \begin{codeframe}{Naslov} ... \end{codeframe}
% Primjer s opcijama: \begin{codeframe}[shrink=15]{Naslov} ... \end{codeframe}
\newenvironment{codeframe}[2][]
  {\begin{frame}[fragile,shrink=5,#1]{#2}}
  {\end{frame}}

% ============================================================

\documentclass[aspectratio=149,11pt]{beamer}
\usetheme{CambridgeUS}

\usepackage{fontspec}              % radi samo sa XeLaTeX/LuaLaTeX
\setsansfont{Latin Modern Sans}    % sans-serif font
\setmonofont{Latin Modern Mono}    % monospace font za kod
\usepackage[croatian]{babel}
\usepackage{csquotes}              % \enquote{} → hrvatski navodnici „..."
\usepackage{lmodern}
\usepackage{booktabs}
\usepackage{listings}
\usepackage{xcolor}
\usepackage{hyperref}

% ---------- Autor i logo ----------
\author{mag.ing. Josip Torić (SIT UNIZD)}
\titlegraphic{\includegraphics[height=2.5cm]{../common/SIT_logo_tekst1.pdf}}

% ---------- Listings ----------
\definecolor{codebg}{RGB}{248,248,248}
\definecolor{codeframeclr}{RGB}{220,220,220}
\lstdefinelanguage{bash}{
  morekeywords={sudo,apt,brew,pip,pipx,uv,winget,python,git,flask,source,deactivate,code,gunicorn,docker,compose,heroku,render,railway,node,npm,pnpm,yarn,uvicorn,alembic,pytest,ruff,black,precommit},
  sensitive=false,
  morecomment=[l]{\#}
}
\lstdefinelanguage{python}{
  morekeywords={from,import,def,class,return,if,elif,else,for,in,while,try,except,raise,with,as,not,and,or,is,lambda,yield,True,False,None,await,async},
  sensitive=true,
  morecomment=[l]{\#},
  morestring=[b]',
  morestring=[b]"
}
\lstdefinelanguage{sql}{
  morekeywords={SELECT,FROM,WHERE,INSERT,INTO,VALUES,UPDATE,SET,DELETE,CREATE,TABLE,PRIMARY,KEY,FOREIGN,REFERENCES,UNIQUE,NOT,NULL,INDEX,ON,ALTER,DROP,CONSTRAINT},
  sensitive=false,
  morecomment=[l]{--}
}
\lstdefinelanguage{json}{
  morecomment=[l]{//},
  morestring=[b]",
  showstringspaces=false
}
\lstdefinelanguage{typescript}{
  morekeywords={import,export,from,const,let,var,function,return,if,else,for,of,in,while,switch,case,break,async,await,try,catch,finally,throw,new,class,this,extends,implements,interface,type,enum,readonly,public,private,protected,static,get,set,true,false,null,undefined,void,string,number,boolean,any,unknown,never,typeof,keyof,as},
  sensitive=true,
  morecomment=[l]{//},
  morecomment=[s]{/*}{*/},
  morestring=[b]',
  morestring=[b]",
  morestring=[b]`
}
\lstalias{javascript}{typescript}
\lstset{
  backgroundcolor=\color{codebg},
  basicstyle=\ttfamily\footnotesize,
  frame=single,
  rulecolor=\color{codeframeclr},
  frameround=tttt,
  breaklines=true,
  tabsize=2,
  showstringspaces=false,
  upquote=true,
  columns=fullflexible
}

% ---------- Helper macros ----------
\newcommand{\codefile}[1]{\texttt{#1}}
\newcommand{\kbd}[1]{\texttt{#1}}

% ---------- Code frame environment ----------
% Automatski dodaje fragile i shrink=5, koristi se umjesto \begin{frame}
% Primjer: \begin{codeframe}{Naslov} ... \end{codeframe}
% Primjer s opcijama: \begin{codeframe}[shrink=15]{Naslov} ... \end{codeframe}
\newenvironment{codeframe}[2][]
  {\begin{frame}[fragile,shrink=5,#1]{#2}}
  {\end{frame}}


\title[Predavanje 4: Autorizacija]{Programiranje za web \\ Predavanje 4: Autorizacija i vlasništvo nad zapisom}
\date{29.10.2026.}

\begin{document}

\frame{\titlepage}

% ------------------------------------------------------------
\begin{frame}{Plan predavanja}
\begin{enumerate}
  \item Što autorizacija odlučuje
  \item Provjera uloge
  \item Vlasništvo nad zapisom
  \item Odgovori i njihove posljedice
\end{enumerate}

\vspace{0.6em}
Na kraju predavanja svaka ruta zna tko smije poslati zahtjev i nad čijim podatcima.
\end{frame}

% ------------------------------------------------------------
\begin{frame}{Ishodi učenja}
Nakon ovog predavanja student može:

\begin{itemize}
  \item razgraničiti provjeru uloge od provjere vlasništva i reći koju vrstu zaobilaženja svaka sprječava
  \item napisati ovisnost koja propušta samo određene uloge
  \item provesti provjeru vlasništva u sloju poslovnih pravila i obrazložiti zašto ne stoji u ruti
  \item razlikovati odbijanje zahtjeva od sužavanja odgovora
  \item obrazložiti izbor između odgovora \texttt{403} i \texttt{404}
\end{itemize}
\end{frame}

% ------------------------------------------------------------
\begin{frame}{Domena primjera}
Sustav prijava klubova na powerlifting natjecanje. Račun ima ulogu \texttt{admin} ili \texttt{club}. Svaki klub ima jedan račun koji pripada samo njemu; organizator ne pripada nijednom.

\vspace{0.4em}
Iz prethodnog predavanja svaki zahtjev nosi token, a ovisnost \texttt{get\_current\_user} iz njega dobiva zapis o korisniku iz baze. Danas se odlučuje što taj korisnik smije.
\end{frame}

% ============================================================
\sekcija{Što autorizacija odlučuje}{Granica između identiteta i prava, uloga i vlasništvo, modeli autorizacije.}

% ------------------------------------------------------------
\begin{frame}{Od identiteta do prava}
Znati tko šalje zahtjev nije isto što i znati smije li ga poslati. Primjer iz domene:

\vspace{0.5em}
\begin{itemize}
  \item Klub pokušava stvoriti drugi klub. Zahtjev je valjan i pošiljatelj je poznat, ali stvaranje klubova nije njegov posao nego organizatorov.
  \item Klub A dohvaća podatke kluba B. Za tu radnju ima pravo, ali samo nad zapisima svojega kluba.
\end{itemize}

\vspace{0.5em}
Prvi slučaj ovisi samo o tome tko je pošiljatelj. Drugi ovisi i o tome na koji se zapis zahtjev odnosi, pa se ne može odlučiti dok se ne zna o kojem je zapisu riječ.

\vspace{0.3em}
Zato autorizacija ima dva dijela, i oni se provode na različitim mjestima.
\end{frame}

% ------------------------------------------------------------
\begin{frame}{Uloga i vlasništvo}
\begin{center}
{\small
\setlength{\tabcolsep}{8pt}
\renewcommand{\arraystretch}{1.4}
\begin{tabular}{@{}lll@{}}
  \toprule
  \textbf{Provjera} & \textbf{Odgovara na} & \textbf{Sprječava} \\
  \midrule
  uloga      & smije li ovakvu radnju  & radnju iznad svojih prava \\
  vlasništvo & smije li nad ovim zapisom & dohvat tuđih podataka \\
  \bottomrule
\end{tabular}
}
\end{center}

\vspace{0.6em}
\textbf{Uloga} je svojstvo korisnika i ne ovisi o zahtjevu: klub nikada ne stvara klubove.

\vspace{0.3em}
\textbf{Vlasništvo} je odnos između korisnika i zapisa: isti korisnik smije mijenjati jedan zapis, a drugi ne, premda je radnja ista.

\vspace{0.4em}
Iz te razlike slijedi i mjesto provjere: uloga se može provjeriti prije nego bude dohvaćen zapis na koji se zahtjev odnosi, vlasništvo tek nakon toga.
\end{frame}

% ------------------------------------------------------------
\begin{frame}{Modeli autorizacije}
\begin{center}
{\small
\setlength{\tabcolsep}{8pt}
\renewcommand{\arraystretch}{1.4}
\begin{tabular}{@{}lll@{}}
  \toprule
  \textbf{Model} & \textbf{Odluka se donosi prema} & \textbf{Primjer} \\
  \midrule
  ACL  & popisu uz svaki zapis        & prava nad datotekom \\
  RBAC & ulozi korisnika              & organizator i klub \\
  ABAC & svojstvima korisnika, zapisa i okolnosti & pravo koje ovisi o dobu dana \\
  \bottomrule
\end{tabular}
}
\end{center}

\vspace{0.5em}
Redom: popis prava uz svaki zapis (engl. \emph{access control list}), prava vezana uz ulogu (engl. \emph{role-based access control}), prava izračunata iz svojstava (engl. \emph{attribute-based access control}).

\vspace{0.5em}
Ovdje se koristi RBAC uz provjeru vlasništva, jer su uloge dvije i pravila se ne mijenjaju bez izmjene koda. ABAC dolazi u obzir kad pravila ovise o okolnostima ili ih treba mijenjati bez nove inačice aplikacije; ACL kad se prava dodjeljuju pojedinačno, po zapisu.
\end{frame}

% ============================================================
\sekcija{Provjera uloge}{Odakle dolazi uloga, kako se provjerava, uporaba u ruti.}

% ------------------------------------------------------------
\begin{frame}{Odakle dolazi uloga}
Uloga mora biti poznata pri svakom zahtjevu, a doći može s dvaju mjesta:

\vspace{0.4em}
\begin{center}
{\small
\setlength{\tabcolsep}{8pt}
\renewcommand{\arraystretch}{1.4}
\begin{tabular}{@{}lll@{}}
  \toprule
  \textbf{Izvor} & \textbf{Kada promjena vrijedi} & \textbf{Cijena} \\
  \midrule
  iz tokena & tek kad token istekne & nema upita \\
  iz baze   & odmah                 & upit po zahtjevu \\
  \bottomrule
\end{tabular}
}
\end{center}

\end{frame}

% ------------------------------------------------------------
\begin{frame}{Token nosi staro stanje}
Token nosi ono što je vrijedilo u trenutku izdavanja. Oduzme li se korisniku pravo, token izdan prije toga i dalje nosi staru ulogu dok ne istekne. Zapis u bazi nosi novu vrijednost odmah, ali za njega treba upit.

\vspace{0.8em}
\begin{konvencija}[Izbor, ne pravilo]
U našem se projektu uloga čita iz baze, jer \texttt{get\_current\_user} zapis o korisniku ionako dohvaća da provjeri je li račun aktivan; upit je time već plaćen. Aplikacije s vrlo velikim brojem zahtjeva odlučuju iz tokena i prihvaćaju da promjena prava kasni do njegova isteka.
\end{konvencija}
\end{frame}

% ------------------------------------------------------------
\begin{frame}[fragile]{Kako provjeravamo ulogu}
\begin{lstlisting}[language=python]
# app/core/deps.py
def require_role(*allowed_roles: str):
    def checker(current_user: User = Depends(get_current_user)) -> User:
        if current_user.role not in allowed_roles:
            raise AppError("forbidden", "Nemate dozvolu za ovu akciju", 403)
        return current_user
    return checker
\end{lstlisting}

\vspace{0.3em}
\texttt{require\_role} nije ovisnost nego funkcija koja je stvara. Poziv \texttt{require\_role("admin")} vraća funkciju \texttt{checker}, a nju ruta navodi u \texttt{Depends}. Tako jedan zapis pokriva sve kombinacije dopuštenih uloga.

\vspace{0.3em}
Unutar sebe \texttt{checker} navodi \texttt{get\_current\_user} kao svoju ovisnost, pa se provjera identiteta izvede prije provjere uloge. Rutama je time dovoljno navesti samo \texttt{require\_role}.
\end{frame}

% ------------------------------------------------------------
\begin{frame}[fragile]{Uporaba u ruti}
\begin{lstlisting}[language=python]
# app/routers/clubs.py
@router.post("/", response_model=ClubResponse, status_code=201)
async def create_club(
    body: ClubCreate,
    db: AsyncSession = Depends(get_db),
    admin: User = Depends(require_role("admin")),
):
    ...
@router.get("/", response_model=list[ClubResponse])
async def list_clubs(
    db: AsyncSession = Depends(get_db),
    user: User = Depends(require_role("admin", "club")),
):
    ...
\end{lstlisting}

\vspace{0.2em}
Iz potpisa rute vidi se tko je smije pozvati: stvaranje kluba samo organizator, popis klubova i organizator i klub.
\end{frame}

% ============================================================
\sekcija{Vlasništvo nad zapisom}{Zašto ne stoji u ruti, provjera u servisu, sužavanje popisa, zaobilaženje kroz oznaku u putanji.}

% ------------------------------------------------------------
\begin{frame}{Zašto vlasništvo ne stoji u ruti}
Provjera uloge gleda samo korisnika, pa se može izvesti prije nego zahtjev dođe do podataka. Provjera vlasništva to ne može, iz dvaju razloga:

\begin{itemize}
  \item \textbf{Zapis mora biti dohvaćen.} Da bi se usporedilo pripada li klub korisniku, klub treba pročitati iz baze. Ruta to ne radi.
  \item \textbf{Pravilo nije isto za svaku radnju.} Klub smije vidjeti svoje podatke, ali ih ne smije mijenjati; organizator smije oboje. To nije svojstvo putanje nego operacije.
\end{itemize}

\vspace{0.5em}
Vlasništvo je time poslovno pravilo, a ne svojstvo HTTP zahtjeva, pa stoji u sloju \codefile{services}, uz ostala pravila.
\end{frame}

% ------------------------------------------------------------
\begin{frame}[fragile]{Provjera u servisu}
\begin{lstlisting}[language=python]
# app/services/club_service.py
async def _get_or_404(db, club_id: int) -> Club:
    club = await club_repo.get_by_id(db, club_id)
    if not club:
        raise AppError("not_found", "Klub nije pronaden", 404)
    return club

async def get_club(db, club_id: int, current_user: User) -> Club:
    club = await _get_or_404(db, club_id)
    if current_user.role == "club" and current_user.club_id != club_id:
        raise AppError("forbidden", "...", 403)
    return club
\end{lstlisting}

\vspace{0.2em}
Redoslijed nije slučajan: prvo se utvrdi postoji li zapis, pa tek onda smije li ga pošiljatelj vidjeti. Uvjet se odnosi samo na ulogu \texttt{club}, jer organizator smije vidjeti svaki klub, pa se za njega ne provjerava ništa.
\end{frame}

% ------------------------------------------------------------
\begin{frame}[fragile]{Sužavanje popisa}
Zahtjev za popisom ne imenuje nijedan zapis, pa nema zapisa za koji bi se moglo reći da je tuđi. Umjesto odbijanja, odgovor se sužava:

\begin{lstlisting}[language=python]
# app/services/club_service.py
async def list_clubs(db, current_user: User) -> list[Club]:
    if current_user.role == "admin":
        return await club_repo.get_all(db)
    club = await club_repo.get_by_id(db, current_user.club_id)
    return [club] if club else []
\end{lstlisting}

\vspace{0.3em}
Ista ruta, isti statusni kod, različit sadržaj. Klub ne dobiva \texttt{403} nego popis od jednog zapisa.

\vspace{0.3em}
To je i razlog zašto popis nikada ne smije dohvatiti sve pa filtrirati u odgovoru: podatci koji ne pripadaju korisniku ne izlaze iz sloja koji ih čita.
\end{frame}

% ------------------------------------------------------------
\begin{frame}[fragile]{Izmjena oznake u adresi}
Klub 12 otvori popis svojih natjecatelja. Preglednik pošalje:

\begin{lstlisting}
GET /clubs/12/lifters
Authorization: Bearer <token kluba 12>
\end{lstlisting}

\vspace{0.3em}
Isti zahtjev može se poslati i izvan sučelja, s jednom promijenjenom znamenkom:

\begin{lstlisting}
GET /clubs/13/lifters
Authorization: Bearer <token kluba 12>
\end{lstlisting}

\vspace{0.3em}
Token je valjan, nije istekao i uloga je \texttt{club}, kakvu ruta i traži. Provjera uloge dakle prolazi. Bez usporedbe kluba iz tokena s klubom iz putanje, odgovor sadrži natjecatelje kluba 13.

\vspace{0.3em}
Takav se propust naziva nesigurno izravno pozivanje na zapis (engl. \emph{insecure direct object reference}, IDOR).
\end{frame}

% ------------------------------------------------------------
\begin{frame}{Što ga ne sprječava}
\begin{itemize}
  \item \textbf{Sučelje koje tu adresu ne nudi.} Zahtjev ne mora doći iz sučelja. Alat naredbenog retka ili razvojna konzola preglednika šalju ga jednako dobro, s istim tokenom.
  \item \textbf{Oznaka koju je teško pogoditi.} Dovoljno ju je jednom vidjeti, a redni brojevi se ionako pogađaju redom. Nasumična oznaka usporava pretragu, ali ne mijenja ishod kad je oznaka poznata.
  \item \textbf{Provjera samo u putanji.} Oznaka kluba može stići i u tijelu zahtjeva, primjerice pri stvaranju natjecatelja. Provjeri li se samo putanja, tijelo ostaje nepokriveno.
\end{itemize}

\vspace{0.5em}
Sprječava ga jedino usporedba pošiljatelja sa zapisom na koji se zahtjev odnosi, pri svakom zahtjevu i nad svakom oznakom koju je zahtjev donio.
\end{frame}

% ============================================================
\sekcija{Odgovori i njihove posljedice}{Razlika između \texttt{401} i \texttt{403}, izbor između \texttt{403} i \texttt{404}, provjere koje se ne oslanjaju jedna na drugu.}

% ------------------------------------------------------------
\begin{frame}{Kada \texttt{401}, a kada \texttt{403}}
\begin{center}
{\small
\setlength{\tabcolsep}{8pt}
\renewcommand{\arraystretch}{1.4}
\begin{tabular}{@{}lll@{}}
  \toprule
  \textbf{Kod} & \textbf{Stanje} & \textbf{Što klijentu preostaje} \\
  \midrule
  \texttt{401} & pošiljatelj nije utvrđen   & prijaviti se ili obnoviti token \\
  \texttt{403} & pošiljatelj je utvrđen     & ništa, zahtjev neće proći \\
  \bottomrule
\end{tabular}
}
\end{center}

\vspace{0.6em}
Razlika nije formalna nego mijenja ponašanje klijenta. Na \texttt{401} klijent posegne za refresh tokenom i ponovi zahtjev. Vrati li poslužitelj \texttt{401} ondje gdje je zapravo riječ o nedostatku prava, klijent obnovi token, ponovi zahtjev, dobije isto, i korisnik ostane bez objašnjenja.

\vspace{0.3em}
U kodu je razlika vidljiva po tome tko podiže pogrešku: \texttt{get\_current\_user} vraća \texttt{401}, a \texttt{require\_role} i provjera vlasništva \texttt{403}.
\end{frame}

% ------------------------------------------------------------
\begin{frame}{Odati postojanje zapisa ili ne}
Klub traži podatke kluba koji postoji, ali nije njegov. Biramo odgovoriti s:

\begin{itemize}
  \item \texttt{403}: zapis postoji, ali nemaš pravo na njega
  \item \texttt{404}: kao da zapisa nema
\end{itemize}

\vspace{0.5em}
Prvi odgovor odaje da zapis postoji. Nizom zahtjeva s različitim oznakama dobiva se popis onoga što u sustavu postoji, bez ijednog dohvaćenog podatka.

\vspace{0.4em}
\begin{konvencija}[Izbor, ne pravilo]
Za podatke kojima samo postojanje nije tajna, \texttt{403} je jasniji jer korisniku kaže što se dogodilo. Ondje gdje je i postojanje zapisa osjetljivo, \texttt{404} odaje manje. U ovom projektu klubovi su javno poznati, pa se vraća \texttt{403}.
\end{konvencija}
\end{frame}

% ------------------------------------------------------------
\begin{frame}{Zaštita u više slojeva}
\begin{center}
{\small
\setlength{\tabcolsep}{8pt}
\renewcommand{\arraystretch}{1.35}
\begin{tabular}{@{}ll@{}}
  \toprule
  \textbf{Mjesto} & \textbf{Što zaustavlja} \\
  \midrule
  sučelje                & radnju koju korisnik ne treba vidjeti \\
  \codefile{routers}     & zahtjev bez tokena i pogrešnu ulogu \\
  \codefile{services}    & zahtjev nad tuđim zapisom \\
  baza                   & zapis koji krši ograničenje \\
  \bottomrule
\end{tabular}
}
\end{center}

\vspace{0.6em}
Ne smijemo se oslanjati na zaštitu preko sučelja: ono što ne prikaže i dalje se može poslati izravno na API.

\vspace{0.3em}
Provjera uloge ne zaustavlja dohvat tuđeg zapisa, a ograničenje u bazi ne zna ništa o tome tko je poslao zahtjev.
\end{frame}

% ------------------------------------------------------------
\begin{frame}{Zahtjevi koji ne bi smjeli proći}
Rutu obično isprobavamo kao korisnik koji na nju ima pravo, i tada sve radi. Izostane li provjera prava, ništa se ne mijenja: zahtjev i dalje prolazi, a odgovor izgleda jednako. Propust se pokaže tek kad isti zahtjev pošaljemo kao korisnik koji na njega nema pravo.

\vspace{0.5em}
Zato za svaku rutu moramo testirati:

\begin{itemize}
  \item bez tokena ili s neispravnim tokenom, uz očekivani \texttt{401}
  \item s tokenom korisnika pogrešne uloge, uz očekivani \texttt{403}
  \item s tokenom korisnika ispravne uloge, ali nad tuđim zapisom, uz očekivani \texttt{403}
\end{itemize}

\vspace{0.5em}
Treći je slučaj najlakše previdjeti, jer je zahtjev u svemu ispravan osim u jednom broju. Taj je propust ujedno i najčešći.
\end{frame}

% ------------------------------------------------------------
\begin{frame}{Sažetak}
\begin{itemize}
  \item Autorizacija odgovara na dva pitanja: smije li korisnik ovakvu radnju i smije li je nad ovim zapisom.
  \item Uloga se provjerava prije dohvata podataka, u ovisnosti uz rutu; vlasništvo tek nakon što je zapis pročitan, u sloju poslovnih pravila.
  \item Uloga koja ulazi u odluku dolazi iz baze, a ne iz tokena, pa promjena prava vrijedi odmah.
  \item Za popise nema što odbiti, pa se odgovor sužava: ista ruta, isti kod, različit sadržaj.
  \item \texttt{401} znači da se treba prijaviti, \texttt{403} da prijava neće pomoći; izbor između \texttt{403} i \texttt{404} odlučuje odaje li se postojanje zapisa.
\end{itemize}
\end{frame}

\end{document}
